\section{Discussion and Interpretation Caveats}
\label{sec:discussion}

The ARFT dataset's main strength is not just its multi-modal measurements, but its explicit distinction between the logged campaign and the analysis-ready subset. As a result, transparency about the currently merged 9 (rather than 12) experiments is essential. The completeness report should therefore be read on two levels: rover logs show what was physically traversed, while the merged \gls{rf}-acoustic dataset shows what is available for channel analysis. This is why EXP001, EXP002, and EXP004 remain relevant to the campaign narrative, despite not contributing processed \gls{rf} and acoustic data to the current merged set.

The second caveat is structural. The \texttt{cycle\_id} axis in the merged \gls{rf} NetCDF is a shared union axis across experiments, not a guarantee that every experiment occupies every cycle index. Validity is encoded by the availability masks, not merely by coordinate presence. Any downstream study that ignores \texttt{csi\_available} or \texttt{position\_available} risks accidentally treating absent data as zero-valued data or assuming that coordinate existence implies measurement existence.

Furthermore, the measurements were performed statically. Therefore, especially the slow-propagating and long reverberating acoustic waves should be handled as static measurements to perform realistic positioning.
